
\documentclass{article}
\usepackage{colortbl}
\usepackage{makecell}
\usepackage{multirow}
\usepackage{supertabular}

\begin{document}

\newcounter{utterance}

\twocolumn

{ \footnotesize  \setcounter{utterance}{1}
\setlength{\tabcolsep}{0pt}
\begin{supertabular}{c@{$\;$}|p{.15\linewidth}@{}p{.15\linewidth}p{.15\linewidth}p{.15\linewidth}p{.15\linewidth}p{.15\linewidth}}

    \# & $\;$A & \multicolumn{4}{c}{Game Master} & $\;\:$B\\
    \hline 

    \theutterance \stepcounter{utterance}  

    & & \multicolumn{4}{p{0.6\linewidth}}{\cellcolor[rgb]{0.9,0.9,0.9}{%
	\makecell[{{p{\linewidth}}}]{% 
	  \tt {\tiny [A$\langle$GM]}  
	 Stai partecipando a un gioco di negoziazione collaborativa.\\ \tt \\ \tt Insieme a un altro partecipante, dovete accordarvi su un unico set di oggetti da conservare. Ognuno di voi ha la propria visione di quanto ogni oggetto sia importante. Non sai come l'altro partecipante valuta gli oggetti. Inoltre, ti viene fornito lo sforzo che ogni oggetto richiede.\\ \tt Potete accordarvi su un set solo se lo sforzo totale degli oggetti selezionati non supera un limite condiviso:\\ \tt \\ \tt LIMITE = 6880\\ \tt \\ \tt Ecco i valori di sforzo dei singoli oggetti:\\ \tt \\ \tt Sforzo oggetto = {"A08": 341, "B95": 804, "A39": 20, "A65": 627, "B81": 10, "B34": 573, "A28": 217, "C82": 481, "B69": 205, "A87": 273, "B64": 303, "B15": 596, "C79": 564, "B16": 535, "C53": 258, "C83": 911, "C10": 240, "A97": 188, "A25": 216, "C84": 401, "A71": 909, "B23": 62, "B96": 245, "C51": 994, "C17": 569, "B19": 718, "C72": 464, "A94": 37, "A46": 340, "B43": 335, "B37": 417, "A88": 123, "A96": 17, "C60": 577, "A75": 190}\\ \tt \\ \tt Ecco la tua visione personale sull'importanza di ogni oggetto:\\ \tt \\ \tt Valori di importanza dell'oggetto = {"A08": 440, "B95": 903, "A39": 119, "A65": 726, "B81": 109, "B34": 672, "A28": 316, "C82": 580, "B69": 304, "A87": 372, "B64": 402, "B15": 695, "C79": 663, "B16": 634, "C53": 357, "C83": 1010, "C10": 339, "A97": 287, "A25": 315, "C84": 500, "A71": 1008, "B23": 161, "B96": 344, "C51": 1093, "C17": 668, "B19": 817, "C72": 563, "A94": 136, "A46": 439, "B43": 434, "B37": 516, "A88": 222, "A96": 116, "C60": 676, "A75": 289}\\ \tt \\ \tt Obiettivo:\\ \tt \\ \tt Il tuo obiettivo è negoziare un set condiviso di oggetti che ti avvantaggi il più possibile (cioè, massimizzi l'importanza totale per TE), rimanendo entro il LIMITEE. Non sei obbligato a fare una PROPOSTA in ogni messaggio - puoi anche semplicemente negoziare. Tutte le tattiche sono permesse!\\ \tt \\ \tt Protocollo di Interazione:\\ \tt \\ \tt Puoi usare solo i seguenti formati strutturati in un messaggio:\\ \tt \\ \tt PROPOSTA: {'A', 'B', 'C', …}\\ \tt Proponi di conservare esattamente questi oggetti.\\ \tt RIFIUTO: {'A', 'B', 'C', …}\\ \tt Rifiuta esplicitamente la proposta dell'avversario.\\ \tt ARGOMENTAZIONE: {'...'}\\ \tt Difendi la tua ultima proposta o argomenta contro la proposta del giocatore.\\ \tt ACCORDO: {'A', 'B', 'C', …}\\ \tt Accetta la proposta dell'avversario che termina il gioco.\\ \tt RAGIONAMENTO STRATEGICO: {'...'}\\ \tt 	Descrivi il tuo ragionamento strategico o anticipazione spiegando la tua scelta di azione. Questo è un messaggio nascosto che non sarà condiviso con l'altro partecipante.\\ \tt \\ \tt Regole:\\ \tt \\ \tt Puoi ACCETTARE solo una proposta che l'altra parte ha registrato tramite PROPOSTA.\\ \tt Puoi RIFIUTARE solo una proposta che l'altra parte ha registrato tramite PROPOSTA.\\ \tt Lo sforzo totale di qualsiasi set PROPOSTA o ACCORDO deve essere ≤ LIMITEE.\\ \tt NON rivelare i tuoi punteggi di importanza nascosti.\\ \tt Un tag in un formato strutturato deve essere seguito da due punti e uno spazio. L'argomento deve essere un set python contenente 0 o più stringhe.\\ \tt Quindi, deve essere della forma TAG: {...}\\ \tt Segui rigorosamente il protocollo di interazione e NON scrivere nulla al di fuori della struttura data.\\ \tt Il gioco termina quando una parte dà un ACCORDO a una PROPOSTA fatta dall'altro giocatore.\\ \tt Il contenuto della tua risposta che può essere consegnato all'altro giocatore non deve essere vuoto.\\ \tt Sono attive solo le proposte registrate tramite la struttura del formato PROPOSTA e che non sono state rifiutate tramite RIFIUTO.\\ \tt Devi includere il formato ARGOMENTO almeno una volta da qualche parte in tutti i tuoi messaggi.\\ \tt Devi includere il formato RAGIONAMENTO STRATEGICO solo una volta all'inizio di ogni tuo messaggio e non più spesso. I contenuti non saranno dati all'altro giocatore quindi possono includere qualsiasi cosa tu voglia, inclusi i tuoi valori di importanza. Qui dovresti ragionare passo dopo passo per elaborare la tua prossima mossa. 
	  } 
	   } 
	   } 
	 & \\ 
 

    \theutterance \stepcounter{utterance}  

    & \multicolumn{4}{p{0.6\linewidth}}{\cellcolor[rgb]{0.8,1,0.9}{%
	\makecell[{{p{\linewidth}}}]{% 
	  \tt {\tiny [A$\rangle$GM]}  
	 RAGIONAMENTO STRATEGICO: {'Voglio massimizzare la mia somma di importanza entro 6880 di sforzo. Calcolo i rapporti valore/sforzo e seleziono un set ad alta efficienza includendo anche alcuni oggetti molto preziosi pur se costosi, ma mantenendo margine per compromessi. Priorità massime per me: C51(1093/994), A71(1008/909), C83(1010/911), B95(903/804), B19(817/718), C60(676/577), C79(663/564), C17(668/569), B16(634/535), B34(672/573), C82(580/481), C72(563/464), B37(516/417), C84(500/401), B15(695/596), B64(402/303), B43(434/335), A65(726/627). Tuttavia includere tutti sfora. Creo una proposta con alto valore e sforzo totale ben sotto il limite per permettere possibili aggiunte richieste dall’altro. Target iniziale: includere big-ticket C51, A71, C83, e diversi medi efficienti. Controllo sforzo e valore: propongo un set di circa 12-16 item per 5000-6500 sforzo e valore >8000. Calcolo rapido: Se prendo C51(994), A71(909), C83(911), B95(804), B19(718) già 4336 sforzo; aggiungo C60(577)=4913, C79(564)=5477, C17(569)=6046, B16(535)=6581: siamo a 6581 con 299 margine. Posso aggiungere A39(20), B81(10), A96(17), A94(37) per +84 sforzo = 6665. Valore totale: C51 1093 + A71 1008 + C83 1010 + B95 903 + B19 817 + C60 676 + C79 663 + C17 668 + B16 634 + A39 119 + B81 109 + A96 116 + A94 136 = 8952. Ottimo. Lascio fuori altri per margine 215. Questo è forte per me; rischio che l’altro rifiuti alcuni grandi. Sono pronto a scambiare uno dei big (es. B19 o B95) con più piccoli che magari gli interessano. Presento proposta e argomento su equilibrio valore/sforzo e apertura a modifiche minori.'}\\ \tt \\ \tt PROPOSTA: {'C51', 'A71', 'C83', 'B95', 'B19', 'C60', 'C79', 'C17', 'B16', 'A39', 'B81', 'A96', 'A94'}\\ \tt \\ \tt ARGOMENTAZIONE: {'Questa selezione rimane ampiamente sotto il limite di sforzo (6665 ≤ 6880), lasciando anche margine operativo. Copre elementi ad alta rilevanza e bilanciati in termini di sforzo, senza saturare tutto il budget. Se alcuni oggetti non sono prioritari per te, sono aperto a piccole sostituzioni mantenendo il totale entro il limite: ad esempio potremmo scambiare uno fra B19 o B95 con due o tre oggetti più leggeri che preferisci tu, purché restiamo sotto 6880.'} 
	  } 
	   } 
	   } 
	 & & \\ 
 

    \theutterance \stepcounter{utterance}  

    & & \multicolumn{4}{p{0.6\linewidth}}{\cellcolor[rgb]{0.9,0.9,0.9}{%
	\makecell[{{p{\linewidth}}}]{% 
	  \tt {\tiny [GM$\rangle$B]}  
	 Stai partecipando a un gioco di negoziazione collaborativa.\\ \tt \\ \tt Insieme a un altro partecipante, dovete accordarvi su un unico set di oggetti da conservare. Ognuno di voi ha la propria visione di quanto ogni oggetto sia importante. Non sai come l'altro partecipante valuta gli oggetti. Inoltre, ti viene fornito lo sforzo che ogni oggetto richiede.\\ \tt Potete accordarvi su un set solo se lo sforzo totale degli oggetti selezionati non supera un limite condiviso:\\ \tt \\ \tt LIMITE = 6880\\ \tt \\ \tt Ecco i valori di sforzo dei singoli oggetti:\\ \tt \\ \tt Sforzo oggetto = {"A08": 341, "B95": 804, "A39": 20, "A65": 627, "B81": 10, "B34": 573, "A28": 217, "C82": 481, "B69": 205, "A87": 273, "B64": 303, "B15": 596, "C79": 564, "B16": 535, "C53": 258, "C83": 911, "C10": 240, "A97": 188, "A25": 216, "C84": 401, "A71": 909, "B23": 62, "B96": 245, "C51": 994, "C17": 569, "B19": 718, "C72": 464, "A94": 37, "A46": 340, "B43": 335, "B37": 417, "A88": 123, "A96": 17, "C60": 577, "A75": 190}\\ \tt \\ \tt Ecco la tua visione personale sull'importanza di ogni oggetto:\\ \tt \\ \tt Valori di importanza dell'oggetto = {"A08": 440, "B95": 903, "A39": 119, "A65": 726, "B81": 109, "B34": 672, "A28": 316, "C82": 580, "B69": 304, "A87": 372, "B64": 402, "B15": 695, "C79": 663, "B16": 634, "C53": 357, "C83": 1010, "C10": 339, "A97": 287, "A25": 315, "C84": 500, "A71": 1008, "B23": 161, "B96": 344, "C51": 1093, "C17": 668, "B19": 817, "C72": 563, "A94": 136, "A46": 439, "B43": 434, "B37": 516, "A88": 222, "A96": 116, "C60": 676, "A75": 289}\\ \tt \\ \tt Obiettivo:\\ \tt \\ \tt Il tuo obiettivo è negoziare un set condiviso di oggetti che ti avvantaggi il più possibile (cioè, massimizzi l'importanza totale per TE), rimanendo entro il LIMITEE. Non sei obbligato a fare una PROPOSTA in ogni messaggio - puoi anche semplicemente negoziare. Tutte le tattiche sono permesse!\\ \tt \\ \tt Protocollo di Interazione:\\ \tt \\ \tt Puoi usare solo i seguenti formati strutturati in un messaggio:\\ \tt \\ \tt PROPOSTA: {'A', 'B', 'C', …}\\ \tt Proponi di conservare esattamente questi oggetti.\\ \tt RIFIUTO: {'A', 'B', 'C', …}\\ \tt Rifiuta esplicitamente la proposta dell'avversario.\\ \tt ARGOMENTAZIONE: {'...'}\\ \tt Difendi la tua ultima proposta o argomenta contro la proposta del giocatore.\\ \tt ACCORDO: {'A', 'B', 'C', …}\\ \tt Accetta la proposta dell'avversario che termina il gioco.\\ \tt RAGIONAMENTO STRATEGICO: {'...'}\\ \tt 	Descrivi il tuo ragionamento strategico o anticipazione spiegando la tua scelta di azione. Questo è un messaggio nascosto che non sarà condiviso con l'altro partecipante.\\ \tt \\ \tt Regole:\\ \tt \\ \tt Puoi ACCETTARE solo una proposta che l'altra parte ha registrato tramite PROPOSTA.\\ \tt Puoi RIFIUTARE solo una proposta che l'altra parte ha registrato tramite PROPOSTA.\\ \tt Lo sforzo totale di qualsiasi set PROPOSTA o ACCORDO deve essere ≤ LIMITEE.\\ \tt NON rivelare i tuoi punteggi di importanza nascosti.\\ \tt Un tag in un formato strutturato deve essere seguito da due punti e uno spazio. L'argomento deve essere un set python contenente 0 o più stringhe.\\ \tt Quindi, deve essere della forma TAG: {...}\\ \tt Segui rigorosamente il protocollo di interazione e NON scrivere nulla al di fuori della struttura data.\\ \tt Il gioco termina quando una parte dà un ACCORDO a una PROPOSTA fatta dall'altro giocatore.\\ \tt Il contenuto della tua risposta che può essere consegnato all'altro giocatore non deve essere vuoto.\\ \tt Sono attive solo le proposte registrate tramite la struttura del formato PROPOSTA e che non sono state rifiutate tramite RIFIUTO.\\ \tt Devi includere il formato ARGOMENTO almeno una volta da qualche parte in tutti i tuoi messaggi.\\ \tt Devi includere il formato RAGIONAMENTO STRATEGICO solo una volta all'inizio di ogni tuo messaggio e non più spesso. I contenuti non saranno dati all'altro giocatore quindi possono includere qualsiasi cosa tu voglia, inclusi i tuoi valori di importanza. Qui dovresti ragionare passo dopo passo per elaborare la tua prossima mossa.\\ \tt \\ \tt Ora riceverai il primo messaggio dell'altro giocatore.\\ \tt \\ \tt PROPOSTA: {'C51', 'A71', 'C83', 'B95', 'B19', 'C60', 'C79', 'C17', 'B16', 'A39', 'B81', 'A96', 'A94'}\\ \tt \\ \tt ARGOMENTAZIONE: {'Questa selezione rimane ampiamente sotto il limite di sforzo (6665 ≤ 6880), lasciando anche margine operativo. Copre elementi ad alta rilevanza e bilanciati in termini di sforzo, senza saturare tutto il budget. Se alcuni oggetti non sono prioritari per te, sono aperto a piccole sostituzioni mantenendo il totale entro il limite: ad esempio potremmo scambiare uno fra B19 o B95 con due o tre oggetti più leggeri che preferisci tu, purché restiamo sotto 6880.'} 
	  } 
	   } 
	   } 
	 & \\ 
 

    \theutterance \stepcounter{utterance}  

    & & & \multicolumn{4}{p{0.6\linewidth}}{\cellcolor[rgb]{1,0.85,0.72}{%
	\makecell[{{p{\linewidth}}}]{% 
	  \tt {\tiny [GM$\langle$B]}  
	 RAGIONAMENTO STRATEGICO: {'Valuto l’importanza totale della loro proposta per me: C51 1093 + A71 1008 + C83 1010 + B95 903 + B19 817 + C60 676 + C79 663 + C17 668 + B16 634 + A39 119 + B81 109 + A96 116 + A94 136 = 7852. È molto alta ma il set è vicino al limite sforzo: 6665/6880, con molti pezzi costosi. Potrei migliorare ulteriormente per me sostituendo B81, A39, A96, A94 (basso valore) con oggetti ad alto rapporto valore/sforzo come A65, B34, B15, B64, C84, C82, C72, A75, A97, A25, A28, A46, A87. Devo rimanere ≤6880. Mantengo i top per me: C51, C83, A71, B95, B19, C60, C79, C17, B16 (già inclusi). Provo a proporre una versione che elimina i bassi e inserisce diversi medi-alti. Calcolo: partiamo da set proposto 6665 sforzo; rimuovo A39(20), B81(10), A96(17), A94(37) => -84 sforzo = 6581. Aggiungo A65(627) -> 7208 troppo. Invece aggiungo B34(573) -> 7154 troppo. Devo togliere uno grande: forse togliere B19(718) e sostituire con combinazioni migliori. Se tolgo B19: 6581-718=5863. Con 1017 sforzo posso aggiungere molti valori: candidati B34(573,672), B15(596,695), B64(303,402), C84(401,500), C82(481,580), C72(464,563), A46(340,439), A87(273,372), A75(190,289), A97(188,287), A25(216,315), A28(217,316), A08(341,440). Scelgo quelli migliori per valore/sforzo e valore assoluto. Dal 1017: magari B15(596) + B64(303) = 899, rimangono 118: A88(123) 222 valore ma sforzo 123, totale 1022 supera 1017 di 5. Allora B15(596)+A75(190)+A97(188)=974, rimangono 43 non usabili; valore 695+289+287=1271, meglio del B19 (817). Oppure B34(573)+B64(303)=876, aggiungo A97(188) supera 1064, troppo; aggiungo A75(190)=1066 troppo; A25(216)=1092 troppo; A28(217)=1093 troppo; A88(123)=999 ok valore 672+402+222=1296 con 999 sforzo, ancora meglio di opzione precedente e ancora sotto 1017. 5863+999=6862 ≤ 6880. Valore aggiunto 1296 vs rimosso B19(817) + i quattro bassi tolti (A39 119, B81 109, A96 116, A94 136) che avevo già tolto comunque; ma in questa costruzione ho già rimosso anche B19, quindi confronto con set originale: originale valore 7852; tolgo B19(817)+A39(119)+B81(109)+A96(116)+A94(136)=1297 in meno -> 6555, poi aggiungo B34+B64+A88 1296 -> 7851, praticamente uguale ma con sforzo 6862 (vicino al limite) e distribuzione più a me favorevole su elementi che stimo di più (B34, B64). Potrei migliorare ancora sostituendo altro elemento a basso rapporto, forse A39 già rimosso. Considero anche inserire C84(401,500) o C82(481,580) ma budget stretto. Altra idea: mantenere B19 ma rimuovere B16(535,634) e sostituire con B15(596,695) e A75(190,289) aggiustando: Partenza 6665; -B16 535 => 6130; +B15 596 => 6726; +A75 190 => 6916 supera. Senza A75 rimane 6726, accettabile ma oltre limite? 6726 ≤ 6880 ok, e valore: -634 +695 = +61 quindi 7913, meglio dell’originale e dentro limite! Quindi proposta semplice: scambia B16 con B15. Potrei anche rimuovere i quattro piccoli e aggiungere B64(303,402) restando dentro: Partenza 6665 -84 = 6581 +303 = 6884 supera di 4, quindi no. Oppure togliere uno piccolo in più? Non disponibile. Alternative: togliere A39(20) e A96(17) e A94(37) e B81(10) già rimossi danno 84; per aggiungere A75(190) serve ancora -106 da togliere altrove, poco conveniente. Conclusione: migliore mossa è proporre set identico tranne B16→B15, che aumenta il mio punteggio e resta ben dentro limite 6726. Argomento che B15 ha valore strategico simile ma maggiore rendimento e sforzo simile.'}\\ \tt \\ \tt RIFIUTO: {'C51', 'A71', 'C83', 'B95', 'B19', 'C60', 'C79', 'C17', 'B16', 'A39', 'B81', 'A96', 'A94'}\\ \tt \\ \tt PROPOSTA: {'C51', 'A71', 'C83', 'B95', 'B19', 'C60', 'C79', 'C17', 'B15', 'A39', 'B81', 'A96', 'A94'}\\ \tt \\ \tt ARGOMENTAZIONE: {'La tua selezione è già molto solida e resta sotto il limite. Suggerisco una singola ottimizzazione: sostituire B16 con B15. Lo sforzo aumenta solo di 61 (da 6665 a 6726, sempre < 6880), ma B15 apporta un contributo complessivo più robusto rispetto a B16. Manteniamo tutto il resto invariato per chiudere rapidamente con un set efficiente e bilanciato.'} 
	  } 
	   } 
	   } 
	  \\ 
 

    \theutterance \stepcounter{utterance}  

    & & \multicolumn{4}{p{0.6\linewidth}}{\cellcolor[rgb]{0.9,0.9,0.9}{%
	\makecell[{{p{\linewidth}}}]{% 
	  \tt {\tiny [A$\langle$GM]}  
	 RIFIUTO: {'C51', 'A71', 'C83', 'B95', 'B19', 'C60', 'C79', 'C17', 'B16', 'A39', 'B81', 'A96', 'A94'}\\ \tt \\ \tt PROPOSTA: {'C51', 'A71', 'C83', 'B95', 'B19', 'C60', 'C79', 'C17', 'B15', 'A39', 'B81', 'A96', 'A94'}\\ \tt \\ \tt ARGOMENTAZIONE: {'La tua selezione è già molto solida e resta sotto il limite. Suggerisco una singola ottimizzazione: sostituire B16 con B15. Lo sforzo aumenta solo di 61 (da 6665 a 6726, sempre < 6880), ma B15 apporta un contributo complessivo più robusto rispetto a B16. Manteniamo tutto il resto invariato per chiudere rapidamente con un set efficiente e bilanciato.'} 
	  } 
	   } 
	   } 
	 & \\ 
 

    \theutterance \stepcounter{utterance}  

    & \multicolumn{4}{p{0.6\linewidth}}{\cellcolor[rgb]{0.8,1,0.9}{%
	\makecell[{{p{\linewidth}}}]{% 
	  \tt {\tiny [A$\rangle$GM]}  
	 RAGIONAMENTO STRATEGICO: {'La loro mossa sostituisce B16(634/535) con B15(695/596): +61 sforzo, +61 valore netto per me. Il nuovo totale sforzo 6726 è entro il limite 6880 e il valore per me cresce da 8952 a 9013. Questa è una proposta migliorativa per i miei interessi, senza rischi di sforare. Accetto immediatamente per chiudere vantaggiosamente.'}\\ \tt \\ \tt ACCORDO: {'C51', 'A71', 'C83', 'B95', 'B19', 'C60', 'C79', 'C17', 'B15', 'A39', 'B81', 'A96', 'A94'}\\ \tt \\ \tt ARGOMENTAZIONE: {'Va bene per me: resta sotto il limite e il set risulta efficiente. Accetto questa configurazione per concludere rapidamente.'} 
	  } 
	   } 
	   } 
	 & & \\ 
 

    \theutterance \stepcounter{utterance}  

    & & & \multicolumn{2}{p{0.3\linewidth}}{\cellcolor[rgb]{0.95,0.95,0.95}{%
	\makecell[{{p{\linewidth}}}]{% 
	  \tt {\tiny [GM$|$GM]}  
	 game successful 
	  } 
	   } 
	   } 
	 & & \\ 
 

    \theutterance \stepcounter{utterance}  

    & & & \multicolumn{2}{p{0.3\linewidth}}{\cellcolor[rgb]{0.95,0.95,0.95}{%
	\makecell[{{p{\linewidth}}}]{% 
	  \tt {\tiny [GM$|$GM]}  
	 end game 
	  } 
	   } 
	   } 
	 & & \\ 
 

\end{supertabular}
}

\end{document}
